% =============================================================================
% CHAPTER 4: Mathematical Framework & Stability
% Sources: part1_mathematical_framework.md, part3 §6.1.3, theory.md (stability),
%          part2 §5.4 (quantum loop corrections)
% =============================================================================

\chapter{Mathematical Framework \& Stability}
\label{ch:mathematical}

% -----------------------------------------------------------------------------
\section{Formal Action Formulation}
\label{sec:action}

The theory extends the Randall--Sundrum framework to include dynamic brane fluctuations:
\begin{equation}
S = \int d^5x\,\sqrt{-g_5}\left[\frac{M_5^3}{2}\,R_5 - \Lambda_5\right]
+ \int d^4x\,\sqrt{-g_4}\left[\frac{M_P^2}{2}\,R_4 - \tau(t,\vec{x}) + \mathcal{L}_\text{matter}\right]
\end{equation}
where:
\begin{itemize}
  \item $M_5$ is the 5D Planck mass,
  \item $\Lambda_5$ is the bulk cosmological constant,
  \item $\tau(t,\vec{x})$ is the dynamic brane tension,
  \item $\mathcal{L}_\text{matter}$ includes all Standard Model fields.
\end{itemize}

% -----------------------------------------------------------------------------
\section{The Radion Field}
\label{sec:radion}

Brane oscillations are described by a scalar field $\phi(x)$ representing the brane's position in the extra dimension:
\begin{equation}
\tau(t,\vec{x}) = \tauZ + \delta\tau\,\cos(\omega t + \vec{k}\cdot\vec{x})
\end{equation}
where oscillations satisfy the Klein--Gordon equation in the bulk:
\begin{equation}
\Box_5\,\phi + m_\phi^2\,\phi = 0
\end{equation}

The effective 4D action after integrating out the extra dimension:
\begin{equation}
S_\text{eff} = \int d^4x\,\sqrt{-g}\left[\frac{M_P^2}{2}\,R + \frac{1}{2}(\partial\phi)^2 - V(\phi) + \phi\,T^\mu{}_\mu\right]
\end{equation}

% -----------------------------------------------------------------------------
\section{Gravitational Effects}
\label{sec:grav_effects}

The oscillating brane induces an effective energy-momentum tensor:
\begin{equation}
T_{\mu\nu}^\text{osc} = \frac{\tauZ\,f_\text{osc}}{M_P^2}\left[g_{\mu\nu} - \frac{1}{2}\,\partial_\mu\phi\,\partial_\nu\phi\right]
\end{equation}

This mimics cold dark matter with:
\begin{itemize}
  \item Zero pressure in the averaged limit,
  \item Energy density $\rho_\text{eff} = \tauZ\,f_\text{osc}/R_H$,
  \item Clustering properties similar to CDM.
\end{itemize}

% -----------------------------------------------------------------------------
\section{Stability Mechanisms}
\label{sec:stability_mechanisms}

To ensure stability and prevent runaway oscillations, we implement a Goldberger--Wise mechanism:
\begin{equation}
V(\phi) = \lambda\left(\phi^2 - v^2\right)^2
\end{equation}

This stabilizes the radion with mass:
\begin{equation}
m_\phi = 2\lambda v \approx \frac{1}{\eV} \times \left(\frac{L}{0.2\,\mu\text{m}}\right)^{-1}
\end{equation}

\textbf{Guaranteed stability}: The predicted maximum amplitude $\delta\tau/\tauZ \sim 10^{-4}$ remains far below the fragmentation threshold ($\delta\tau/\tauZ > 1$). The membrane can oscillate eternally without risk of tearing.

Higher modes ($\ell \ge 2$) have frequencies:
\begin{equation}
\omega_\ell \simeq \sqrt{\ell(\ell+1)} \times \omega_0
\end{equation}
For $\ell = 2$, the frequency is already $\sqrt{6} \approx 2.5$ times higher. Since the source $\Pi(t)$ is quasi-monochromatic at $\omega_0$, coupling to higher modes decreases as $\delta\omega^{-2}$, naturally damping them.

% -----------------------------------------------------------------------------
\section{Compatibility with General Relativity --- Solar System Tests}
\label{sec:gr_compat}

The model must reproduce all GR successes. We ensure this by environmental suppression of the oscillation amplitude:
\begin{equation}
A_\text{osc}(r) = A_0\exp\!\left(-\frac{\rho_\text{local}}{\rho_\text{crit}}\right)
\end{equation}
where $\rho_\text{crit} \sim 10^{-26}$~kg/m$^3$ (galactic density scale). This ensures negligible effects in the Solar System ($\rho \gg \rho_\text{crit}$).

\paragraph{Mercury Perihelion Precession.}
The additional precession from brane oscillations:
\begin{equation}
\delta\dot{\omega} = \frac{3n}{2}\,\frac{A_\text{osc}^2\,\omega_0^2\,r_\text{Merc}^2}{c^2}\,\sin(2\omega_0 t)
\end{equation}
where $n$ is Mercury's mean motion. For Solar System density: $A_\text{osc}(\text{Solar System}) < 10^{-12}$, yielding $\delta\dot{\omega} < 0.01$~arcsec/century compared to GR's 42.98~arcsec/century (observed: $42.98 \pm 0.04$).

\paragraph{Light Deflection.}
The oscillation contribution to deflection angle:
\begin{equation}
\delta\alpha = \frac{4G\Msun}{c^2\,b} \times \frac{A_\text{osc}^2}{2} < 10^{-9}\,\alpha_\text{GR}
\end{equation}
where $b$ is the impact parameter and $\alpha_\text{GR} = 1.75$~arcsec for grazing rays.

\paragraph{Gravitational Redshift.} Unaffected as the time-averaged metric remains unchanged.

\paragraph{Fifth Force Constraints.}
Any scalar-mediated force is suppressed by:
\begin{equation}
\alpha = \frac{\phi\,M_P}{M_5^2} < 10^{-5}
\end{equation}
satisfying E\"ot-Wash experiments.

% -----------------------------------------------------------------------------
\section{Quantum Regime}
\label{sec:quantum_regime}

\paragraph{Particle Content.}
Oscillation quanta (branons) have:
\begin{itemize}
  \item Mass: $m_\text{branon} \sim 1\eV$,
  \item Coupling to SM: gravitational only,
  \item Production rate: negligible at collider energies.
\end{itemize}

\paragraph{Quantum Stability.}
The effective potential prevents cascading:
\begin{equation}
\Gamma_\text{decay} \sim \frac{m_\phi^5}{M_5^6} < H_0
\end{equation}
ensuring cosmological stability.

\paragraph{Loop Corrections.}
One-loop corrections to the brane tension:
\begin{equation}
\delta\tau_\text{1-loop} = \frac{N_\text{KK}\,m_\text{KK}^4}{64\pi^2}\,\ln\!\left(\frac{\Lambda_\text{UV}}{m_\text{KK}}\right)
\end{equation}
remain small for $\Lambda_\text{UV} \lesssim M_5$.

% -----------------------------------------------------------------------------
\section{Observational Confrontations}
\label{sec:obs_confrontations}

\subsection{CMB Anisotropies (Planck Constraints)}

The effective dark matter density at recombination: $\Omega_\text{osc}(z_\text{rec}) = \Omega_\text{CDM} = 0.258 \pm 0.011$. Modifications to the standard $C_\ell$: $\Delta C_\ell / C_\ell < 10^{-3}$ for $\ell < 2000$, achieved by ensuring adiabatic initial conditions. Spectral index unmodified: $n_s = 0.9649 \pm 0.0042$.

\subsection{Galaxy Rotation Curves}

The brane oscillation creates an effective potential:
\begin{equation}
\Phi_\text{eff}(r) = \Phi_\text{baryon}(r) + \Phi_\text{osc}(r)
\end{equation}
where:
\begin{equation}
\Phi_\text{osc}(r) = -\frac{G\,M_\text{osc}}{r}\left[1 - \exp\!\left(-\frac{r}{r_s}\right)\right]
\end{equation}
with scale radius $r_s \sim 10$~kpc, naturally explaining flat rotation curves.

\textbf{Tully--Fisher Relation:} The model predicts $v_\text{flat}^4 = G\,M_\text{baryon}\,a_0$ with $a_0 = cH_0/(2\pi) \times 1.05 = 1.1 \times 10^{-10}$~m/s$^2$.

\subsection{Gravitational Lensing and the Bullet Cluster}

The Bullet Cluster (1E~0657-56) provides a crucial test. During collision:
\begin{enumerate}
  \item \textbf{Initial State}: Two clusters approaching with relative velocity $\sim 4700\kms$.
  \item \textbf{During Collision}: Gas experiences ram pressure; oscillation field passes through unimpeded (no self-interaction).
  \item \textbf{Post-Collision} ($t > 100$~Myr): Gas lags behind by $\Delta x \sim 150$~kpc; oscillation field remains centered on galaxies.
  \item \textbf{Observational Signature}: $\kappa_\text{lensing}(x) = \kappa_\text{galaxies}(x) + \kappa_\text{osc}(x) \neq \kappa_\text{gas}(x)$.
\end{enumerate}

The mass centroid from weak lensing follows the oscillation field (centered on galaxies), while X-ray emission traces the shocked gas---exactly as observed. This provides a natural explanation without particle dark matter.

\subsection{ISW Effect Signature}

The 2~Gyr oscillation creates a unique CMB signature through the Integrated Sachs--Wolfe effect:
\begin{equation}
\frac{\Delta T}{T} \propto \int \dot{\Phi}(t)\,dt
\end{equation}
with oscillating $w(z)$ causing gravitational potentials to pulsate. Resonance peak at $\ell = 10$--$20$ with $\chi^2$ improvement of 32.9 ($6\sigma$ over \LCDM).

% -----------------------------------------------------------------------------
\section{The Adiabatic Shield}
\label{sec:adiabatic}

The brane oscillation frequency is $\nu \sim 1.6 \times 10^{-17}$~Hz (period $2\Gyr$), while the lightest Kaluza--Klein excitations have mass $\sim 1\eV$, corresponding to $\nu_\text{KK} \sim 10^{14}$~Hz. The ratio is:
\begin{equation}
\frac{\nu_\text{brane}}{\nu_\text{KK}} \sim 10^{-31}
\end{equation}

Particle creation is suppressed by a Schwinger factor:
\begin{equation}
\Gamma_\text{branon} \propto e^{-\pi m_\text{KK}^2/(eE)} \sim e^{-10^{31}} \approx 0
\end{equation}

The cosmic yoyo oscillates without quantum friction.

% -----------------------------------------------------------------------------
\section{Quantum Loop Corrections in Curved Background}
\label{sec:quantum_loops}

Quantum corrections in the warped geometry present unique challenges beyond flat-space field theory. The curved background modifies vacuum fluctuations, leading to several important effects.

\subsection{Casimir Energy in Warped Geometry}

The Casimir energy density between two branes separated by distance $L$ in AdS$_5$ (Flachi \& Tanaka 2003):
\begin{equation}
\rho_\text{Casimir}(z) = -\frac{\pi^2}{1440}\,\frac{N_\text{fields}}{z^4}\left[1 + \frac{45}{2\pi^2}\,\zeta(3)\,e^{-2kz} + \mathcal{O}(e^{-4kz})\right]
\end{equation}
where $N_\text{fields}$ is the total degrees of freedom (SM: $\sim 100$), $k$ is the AdS curvature scale, and $\zeta(3) \approx 1.202$.

For oscillating branes, this creates a time-dependent contribution:
\begin{equation}
V_\text{Casimir}(t) = V_0 + V_1\cos(2\omega_0 t) + V_2\cos(4\omega_0 t) + \cdots
\end{equation}
leading to frequency shift $\delta\omega/\omega_0 \sim 10^{-4}(N_\text{fields}/100)$.

\subsection{One-Loop Effective Action}

The one-loop correction from bulk gravitons and matter fields (Garriga, Pujol\`as \& Tanaka 2001):
\begin{equation}
\Gamma_\text{1-loop} = \frac{1}{2}\,\text{Tr}\ln\!\left[-\Box + m^2 + \xi R\right]
\end{equation}

After regularization and renormalization:
\begin{equation}
V_\text{eff}(z) = V_\text{tree}(z) + \frac{1}{64\pi^2}\sum_i (-1)^{F_i}\,n_i\,m_i^4(z)\,\ln\!\left(\frac{m_i^2(z)}{\mu^2}\right)
\end{equation}
where $F_i$ is the fermion number, $n_i$ the degrees of freedom, $m_i(z)$ the field-dependent masses, and $\mu$ the renormalization scale.

For the radion specifically:
\begin{equation}
V_\text{radion}^\text{1-loop} = \frac{3k^4}{32\pi^2}\,z^4\left[\ln(kz) - \frac{1}{4}\right] + \text{counterterms}
\end{equation}

\subsection{Radion Quantization and Stability}

The quantized radion field has peculiar properties due to the warped geometry (Csaki et~al.\ 2000). The mass spectrum:
\begin{equation}
m_n^2 = \frac{4k^2}{9}\left[4 + n(n+3)\right]e^{-2kL}
\end{equation}

For $n = 0$ (radion): $m_\text{radion} = \frac{4k}{3}\,e^{-kL} \approx 1\eV$.

Quantum stability conditions:
\begin{enumerate}
  \item Coleman--Weinberg potential must be bounded below.
  \item Decay rate: $\Gamma_{\text{radion} \to 2\gamma} < H_0$.
  \item Vacuum stability: $\langle\delta z^2\rangle < L^2$.
\end{enumerate}

\subsection{Dynamic Casimir Effect During Oscillations}

The oscillating brane creates particles from vacuum. Particle creation rate (Brevik et~al.\ 2003):
\begin{equation}
\frac{dN}{dt} = \frac{A_\text{brane}}{(2\pi)^3}\int d^3k\,|\beta_k|^2\,\omega_k
\end{equation}
where $\beta_k$ are Bogoliubov coefficients satisfying:
\begin{equation}
|\beta_k|^2 = \frac{\omega_0^2\,A_\text{osc}^2}{4\omega_k^2}\,\sinh^2\!\left(\frac{\pi\omega_k}{aH}\right)
\end{equation}

This leads to:
\begin{itemize}
  \item Energy dissipation: $\dot{E}/E \sim 10^{-5}\,H_0$ (negligible).
  \item Particle spectrum: Thermal with $T_\text{eff} \sim \hbar\omega_0$.
  \item Backreaction: Modifies equation of state by $\Delta w \sim 10^{-6}$.
\end{itemize}

\subsection{Loop Corrections to Israel Junction Conditions}

At one-loop, the junction conditions receive corrections:
\begin{equation}
[K_{\mu\nu}] = -\kappa_5^2\left(T_{\mu\nu} - \frac{1}{3}\,g_{\mu\nu}\,T + T_{\mu\nu}^\text{quantum}\right)
\end{equation}
where:
\begin{equation}
T_{\mu\nu}^\text{quantum} = \frac{1}{16\pi^2}\sum_i n_i\,\langle T_{\mu\nu}^{(i)}\rangle_\text{ren}
\end{equation}

This modifies:
\begin{itemize}
  \item Brane tension renormalization: $\tau_\text{ren} = \tauZ + \delta\tau_\text{quantum}$.
  \item Induced cosmological constant: $\Lambda_\text{ind} = \Lambda_0 + \frac{\pi^2 N}{1440 L^4}$.
  \item Effective Newton's constant: $\Geff = \GN(1 + \alpha\ln(r/L))$.
\end{itemize}

% -----------------------------------------------------------------------------
\section{Energy of the Membrane}
\label{sec:membrane_energy}

The deformation energy of the cosmic membrane is:
\begin{equation}
E_\text{tens} = \frac{1}{2}\,\tauZ\,A\left(\frac{2\pi z}{\lambda}\right)^2
\end{equation}
where $\tauZ = 7.0 \times 10^{19}$~J/m$^2$ is the brane tension, $A \simeq R_H^2$ is the area of the observable universe, $z$ is the displacement in the extra dimension, and $\lambda \simeq 2R_H$ is the fundamental wavelength.
